\documentclass[12pt]{article}
\PassOptionsToPackage{table}{xcolor}
\usepackage[utf8]{inputenc}
\usepackage[T1]{fontenc}
\usepackage[italian]{babel}
\usepackage{multicol}
\usepackage{enumitem}
\usepackage{array}
\usepackage{longtable}
\usepackage{booktabs}
\usepackage{xcolor}
\usepackage{graphicx}
\usepackage{tikz}
\usetikzlibrary{
    shapes.geometric,
    arrows.meta,
    positioning,
    calc,
    patterns,
    decorations.pathmorphing
}
\usepackage[margin=2cm]{geometry}
\usepackage{amsmath,amssymb}
\setlength{\parindent}{0pt}
\definecolor{headercolor}{RGB}{220,220,220}
\newcounter{quesito}
\setcounter{quesito}{0}
\newcommand{\nuovoquesito}{\stepcounter{quesito}\textbf{Domanda \arabic{quesito}.}}
\newcommand{\itemdomanda}[2]{%
    \par\noindent
    \begin{minipage}[t]{\linewidth}
        \nuovoquesito \\ #1
        \begin{enumerate}[label=(\Alph*), nosep, leftmargin=*]
            #2
        \end{enumerate}
    \end{minipage}
    \vspace{10pt}
}
\begin{document}
\section*{Simulazione}
\hfill Data: 17 apr 2026

\itemdomanda{Se A è più alto di B e B è più alto di C, allora:}{
    \item C è più alto di A
    \item A è più basso di C
    \item A è più alto di C
    \item B è il più alto di tutti
    \item Non si può stabilire nulla
}

\itemdomanda{Date le seguenti premesse:
    \begin{enumerate}
        \item Tutti i colibrì hanno un piumaggio coloratissimo.
        \item Nessun uccello grosso si nutre di miele.
        \item Gli uccelli che non si nutrono di miele hanno colori smorti.
    \end{enumerate}
    Ne consegue che:}{
    \item I colibrì sono uccelli piccoli
    \item I colibrì sono uccelli grossi
    \item Gli uccelli con colori smorti sono colibrì
    \item Nessun uccello piccolo si nutre di miele
    \item Tutti gli uccelli grossi hanno il piumaggio colorato
}

\itemdomanda{Date le premesse:
    \begin{enumerate}
        \item Uno squalo non dubita mai di essere ben difeso.
        \item Un pesce che non sappia danzare il minuetto è degno di disprezzo.
        \item Un pesce non è sicuro di essere ben difeso se non ha tre file di denti.
        \item Tutti i pesci, a parte gli squali, sono gentili con i bambini.
        \item Nessun pesce grosso può danzare il minuetto.
        \item Un pesce con tre file di denti non si può disprezzare.
    \end{enumerate}
    Quale conclusione se ne trae?}{
    \item Tutti i pesci grossi sono gentili con i bambini
    \item Gli squali sanno danzare il minuetto
    \item Nessun pesce grosso ha tre file di denti
    \item I pesci gentili con i bambini sono tutti squali
    \item Chi non sa danzare il minuetto è uno squalo
}

\newpage
\section*{Appendice: risposte corrette}
\begin{longtable}{@{}>{\raggedright\arraybackslash}p{0.18\linewidth}>{\raggedright\arraybackslash}p{0.24\linewidth}>{\raggedright\arraybackslash}p{0.46\linewidth}c@{}}
\toprule
\textbf{Materia} & \textbf{Argomento} & \textbf{Titolo} & \textbf{Risp.} \\
\midrule
\endfirsthead
\toprule
\textbf{Materia} & \textbf{Argomento} & \textbf{Titolo} & \textbf{Risp.} \\
\midrule
\endhead
\midrule
\multicolumn{4}{r}{\footnotesize\textit{Continua}} \\
\endfoot
\bottomrule
\endlastfoot
Logica & Logica Deduttiva & Relazione transitiva & C \\
Logica & Logica Deduttiva & Sillogismo dei colibrì & A \\
Logica & Logica Deduttiva & Sillogismo dei pesci grossi & A
\end{longtable}

\end{document}
